\chapter{Literature Review}\label{chap:background}


\section{Emotion in Learning}\label{sec:emotion_in_learning}

Early work in educational psychology treated emotion in achievement settings mainly through the lens of test anxiety. Foundational studies linked anxiety to interference with learning and performance in evaluative situations \cite{mandler1952anxiety}, and later reviews documented substantial correlational and experimental evidence on how anxiety relates to cognition and academic performance \cite{hembree1988testanxiety,zeidner1998testanxiety}. Even within that limited focus, researchers were already trying to understand achievement-related emotions through the way students explain success and failure \cite{weiner1985attributional}. What remained comparatively open, however, was whether anxiety was the only emotion that mattered for how students actually learn in classrooms and during study.

By the early 2000s, that gap was addressed directly. Qualitative and questionnaire-based work showed that students report a wide range of emotions tied to instruction, studying, and testing, not only fear around exams \cite{pekrun_academic_nodate}. The same line of work introduced the term \emph{academic emotions} for this wider group and classified them using two simple dimensions: whether the emotion is positive or negative, and whether it increases or reduces a student's level of activity. Emotions were understood to influence learning and performance indirectly through motivation, learning strategies, mental resources, and self-regulation, while students' results and feedback could in turn influence how they feel in future learning situations \cite{pekrun_academic_nodate}. In other words, researchers increasingly came to understand that learning is affected not only by one negative emotion, but by a broader set of emotional states that may appear during the same lesson and shape the way students pay attention, stay involved, participate, and make progress.

A parallel line of research developed around the concept of engagement rather than focusing only on specific emotion labels, although it still included a clear emotional dimension. In this view, engagement was not treated as a single construct, but as a combination of three related aspects: behavioural participation in learning activities, emotional reactions to school and instruction, and cognitive effort directed towards understanding difficult material \cite{fredricks_httpwwwjstororg_2004}. Emotional engagement in that sense covers interest, boredom, anxiety, and related feelings towards school and academic work, and it is treated as part of why students are willing to keep trying when tasks are hard \cite{fredricks_httpwwwjstororg_2004}. Fredricks and colleagues connected these dimensions of engagement to long-standing motivational theories which suggest that students are more likely to remain involved in learning when they feel competent, have a sense of control over their learning, and feel socially connected to others \cite{connell1991competence}. They also linked them to expectancy-value perspectives, which suggest that students' decisions to engage are shaped by how interesting they find a task and how important or useful they believe it to be \cite{eccles1983expectations}. Engagement therefore became an additional way of understanding affect in education. It did not replace the study of specific emotions, but worked alongside it by offering a broader view of students' emotional involvement in learning.

From the mid-2000s onward, these educational perspectives were strengthened by more explicit theories explaining how emotions arise in learning situations, as well as by neuroscientific work on why affect matters for learning in the first place. One important example is Pekrun's control-value theory of achievement emotions, which argues that students' emotions depend largely on how much control they believe they have over academic tasks and how much value they attach to those tasks. According to this view, different combinations of control and value can give rise to emotions such as enjoyment, hope, anxiety, boredom, and shame, and these emotions can then influence students' interest, motivation, engagement, and academic achievement \cite{pekrun_control-value_2006}. Independently, Immordino-Yang and Damasio argued that higher-level learning in school is not separate from emotion. They described emotional thought as the connection between feelings and the attention, memory, and reasoning processes that education tries to strengthen, and explained that emotion helps make experiences meaningful, which allows knowledge to be remembered and later used in real decisions \cite{immordino-yang_we_2007}. This argument builds on earlier work which suggested that emotion is an essential part of rational behaviour and decision-making, rather than something that simply distracts from clear thinking \cite{damasio1994descartes}. Together, the appraisal perspective and the neuroscience perspective support the same general idea: emotion is not just an extra layer added to learning, but part of the process itself. In this way, these lines of work helped move educational research beyond a purely cognitive view of learning towards one that also recognises the role of motivation and affect.

Research on dialogue-based tutoring then helped explain more clearly how emotions appear during learning, especially when students interact with a tutor while solving problems, asking questions, and receiving feedback. Studies of long tutoring conversations showed that deep learning is often accompanied by a range of emotions, especially boredom, engagement, confusion, and frustration, which appear more often than the traditional basic emotion categories would suggest \cite{dmello_language_2012}. These four states later became the de facto target set in webcam-based educational facial expression recognition benchmarks such as DAiSEE, which the following sections review. Experimental studies showed that a moderate level of confusion can sometimes improve learning outcomes, partly because it increases attention and encourages deeper processing of the material \cite{dmello2014confusionbeneficial}. A later review by Tyng and colleagues connected these findings from classrooms with brain imaging studies that showed how emotion affects thinking and learning \cite{tyng_influences_2017}. Information with emotional meaning usually attracts attention more quickly and is given more focus than neutral information \cite{schupp2007selective,vuilleumier2005brainsbeware}. At the brain level, areas such as the amygdala, medial temporal lobe, and prefrontal cortex help explain how emotional arousal influences the formation and strengthening of memories, which is why emotional events are often remembered more clearly than neutral ones \cite{tyng_influences_2017,sharot2004emotionremembering}. The same review stresses that stress and arousal can help or hurt learning depending on intensity and timing, and that motivational states such as curiosity support exploration and readiness to learn.

More recent review studies still show that emotion is a central part of how students experience learning, especially in difficult subjects and in online settings where students often receive less informal support. A systematic review of emotions in STEM online learning found that emotions strongly influence engagement, motivation, and satisfaction, and that teaching needs a better understanding of students' emotional experiences in order to respond to what they are actually facing \cite{anwar_emotions_2023}. This conclusion builds on a long tradition of educational psychology research that had already treated emotion as an important topic to understand and measure in schools \cite{schutz2007emotioneducation}. Once emotion is treated as part of how students attend, remember, and stay engaged, the practical question becomes how to tell what a learner is feeling while they are actually working. The next section therefore turns from the conceptual case for emotion in learning to the ways researchers have tried to observe or measure it in educational settings.


\section{Emotion Detection in Educational Systems}\label{sec:emotion_detection_in_educational_systems}

Once emotions were established as central to learning, motivation, and achievement, the field faced a more operational question: how can those emotions be detected while learning is actually taking place. Early work in affective computing had already described emotion detection as a measurement problem that could use several different signals rather than only one method \cite{Picard1997AffectiveComputing}, and educational researchers gradually translated that idea into classroom and tutoring contexts \cite{Calvo2010AffectDetection}. What came next was not one simple shift from one method to another, but several overlapping approaches, each with different strengths and limitations in terms of timing, scalability, and practical use in real educational settings.

At first, most educational studies relied on self-report. The achievement-emotions line of work around the \ac{AEQ} gave researchers a structured way to measure recurring states such as enjoyment, pride, anxiety, anger, shame, and boredom in relation to classroom activity and assessment \cite{pekrun_control-value_2006,pekrun_academic_nodate}. The same type of questionnaires is still used in recent classroom AI studies. For example, Omar and colleagues used \ac{AEQ-M} surveys before and after the intervention, along with interviews, to track emotional changes in grade 8 algebra classes using an AI-supported learning environment \cite{omar_academic_2025}. Studies in distance learning also refined self-report methods instead of abandoning them. For example, some studies compared daily questionnaires completed after learning activities with emoticon tools used during the course, where students could mark moments when they felt certain emotions. These methods had clear advantages, including low instrumentation demands, direct access to students' own experiences, and easy deployment across different settings. However, they also had persistent limitations, such as dependence on students' self-awareness, delayed reporting, sensitivity to wording and social context, and difficulty capturing rapid emotional changes in real time \cite{Lavoue2017EmotionalDataCollection}.

Alongside these questionnaire-based approaches, researchers also used direct human observation to understand students' affect and engagement from their visible behaviour in real classrooms. In school psychology, structured coding systems such as \ac{AET-SSBD}, \ac{BOSS}, \ac{DOF}, \ac{COC}, \ac{SECOS}, \ac{ADHD-SOC}, and \ac{SOS} were used to observe students at set time intervals and record whether they were focused, distracted, or showing signs of difficulty \cite{Volpe2005ObservingStudents,Shapiro2004AcademicSkillsProblemsWorkbook}. Related work in learning technologies also used a similar observation-based approach in educational software and classroom interactions \cite{Baker2010BetterFrustratedThanBored,Liu2017DetectStudentsAcademicEmotions}. Compared with questionnaires completed after the event, observation methods gave more detailed information about when emotions and behaviours appeared and linked interpretations more closely to what students were actually doing in the learning setting. However, these methods were demanding because they required observer training, checks for consistency between observers, repeated observations over time, and careful control of problems such as observer drift and reactivity, while still showing only small parts of students' experience during the full course \cite{Volpe2005ObservingStudents}.

As digital tutoring systems collected detailed data about what students did, researchers began to move from manual coding to large-scale behavioural inference. Instead of asking students about their feelings or observing them all the time, they used records such as answer correctness, hint use, timing, and help-seeking patterns to infer affect and disengagement \cite{Baker2007OffTaskBehaviorITS,Cocea2009OffTaskGamingLearning}. Pardos and colleagues are a good example of this shift. They built detectors in \ac{ASSISTments} that did not need sensors and were used to identify states such as boredom, engaged concentration, confusion, frustration, off-task behaviour, and gaming the system. They then followed these affective patterns over a full year and examined how they were related to students' end-of-year mathematics test results \cite{Pardos2014AffectiveStatesStateTests}. Studies in technology-rich classrooms also used coded behaviour data, sequence analysis, and Hidden Markov Models to study how students' emotions and behaviours changed over time \cite{Liu2017DetectStudentsAcademicEmotions}. These approaches made it easier to study many students and to follow them over time, but they still estimated emotion indirectly from behaviour. This is a limitation because the same behaviour can reflect different emotions, the models depend on labelled training data, and their performance may vary across platforms, tasks, and teaching contexts \cite{Pardos2014AffectiveStatesStateTests}. The same behavioural traces later became the default input to many adaptive and reinforcement-learning tutors, even when facial or other perceptual signals were available \cite{zerkouk_comprehensive_2025,meng_self-evolving_2025}.

Dialogue-based tutoring introduced another approach by using language and conversation as signals of emotion. Research on tutoring dialogue showed that emotional states could, in principle, be inferred from what learners and tutors said, even when explicit emotion words were not used \cite{DMello2008AutomaticDetectionLearnersAffect}. In typed \ac{AutoTutor} dialogues, D'Mello and Graesser demonstrated that discourse-sensitive features derived from \ac{LIWC} and \ac{Coh-Metrix} could predict distributions of boredom, confusion, flow/engagement, and frustration over whole sessions \cite{dmello_language_2012,Graesser2004AutoTutorDialogue}. Those results showed that language carries affective information, but only under conditions that are difficult to reproduce in ordinary online study. The models depend on extended tutoring conversation, domain-specific vocabulary, and session-level transcript summaries rather than on the short, ambiguous messages that dominate most digital exercises. A student who types ``I don't get it'', ``ok'', or nothing at all may be confused, bored, or simply unwilling to write; the same string of words can therefore support more than one emotional reading, which makes text a inherently ambiguous channel for real-time adaptation \cite{dmello_language_2012,Lavoue2017EmotionalDataCollection}. Recent chat-based and large-language-model tutors inherit the same problem: they can process what the learner chooses to write, but they cannot see the frustration or disengagement that appears on the face when the learner stops typing \cite{gutierrez_development_2025,meng_self-evolving_2025}. Surveys of sentiment analysis in STEM online learning report the same weakness at scale: text-based methods struggle with sparse emotional language, context dependence, and the gap between what students write and what they actually feel during a task \cite{anwar_emotions_2023,florestiyanto_emotion_2024}.

Facial expression is a stronger and more direct alternative for educational affect monitoring. Comparative work during classroom use of an educational game found that face-only affect detectors were more accurate than detectors built from interaction logs alone, even though facial models could not be applied in every student case \cite{bosch_2015}. That pattern is consistent with a wider lesson from the field: visible affect on the face is closer to the emotional event than inferred labels derived from clicks, timings, or typed responses. Teachers already rely on this signal in classrooms; once webcams became standard in online learning, the same cue could be captured continuously without asking students to describe their feelings in words. Even early affective tutoring systems that experimented with text semantics treated facial expression as a necessary companion rather than an optional extra, because text alone did not provide a stable enough reading of the user's state \cite{lin_employing_2012}. For adaptive systems that must respond while a student is still working, facial behaviour therefore offers a clearer, more continuous, and less ambiguous basis for detecting academic emotions such as confusion, boredom, frustration, and engagement than dialogue or log data alone \cite{lek_academic_2023,vistorte_integrating_2024,petrovica_emotion_2016}. The next section reviews how educational facial expression recognition has developed around that advantage.

\section{Facial Expression Recognition in Education}\label{sec:fer_in_education}

Among the signals reviewed in the previous section, one stood out as visible and easy to capture during everyday learning. A learner's face is what teachers already read in the classroom, and an ordinary webcam is enough to record it during online sessions. This made \ac{FER} an attractive option once deep learning became stable enough to be used outside the laboratory. Over the last decade, a dedicated line of educational FER work has grown around engagement, attention, boredom, confusion, frustration, and related academic emotions, and a systematic review by Lek and Teo confirms that this line is now the dominant computer-vision approach to emotion monitoring in digital learning environments \cite{lek_academic_2023}. More recent syntheses also describe facial behaviour as one of the most practical signals for studying affect in classrooms and online courses \cite{vistorte_integrating_2024,florestiyanto_emotion_2024}.

A central piece of this development was the introduction of a public benchmark for in-the-wild educational affect. Gupta and colleagues released DAiSEE, a dataset of 9{,}068 short webcam clips collected from 112 students, where each clip is labelled with four affective states, namely engagement, boredom, confusion and frustration, and each state is rated on four ordinal levels \cite{gupta_daisee_2016}. Their benchmark also reported the first deep-learning baselines on the dataset, using InceptionNet at frame and video level, a 3D convolutional network (C3D), and a long-term recurrent convolutional network (LRCN). Top-1 accuracy was clearly higher for confusion and frustration than for engagement and boredom, with LRCN reaching 53.7\% on boredom, 57.9\% on engagement, 72.3\% on confusion and 73.5\% on frustration, while frame- and video-level InceptionNet stayed below 50\% on boredom and engagement \cite{gupta_daisee_2016}. Two consequences shaped almost everything that came after. First, DAiSEE became the reference benchmark for academic FER \cite{lek_academic_2023}. Second, the dataset is genuinely difficult, especially for engagement and boredom, and is heavily class-imbalanced, which made temporal modelling, attention and imbalance-aware training the natural next steps. It is also worth noting that within the searched sources, no pure traditional machine-learning paper using DAiSEE as the main method was found; classical components such as facial landmarks, action units or handcrafted features only appear inside hybrid pipelines combined with deep networks.

Early educational FER work in the same period still kept one foot in classical computer vision. Rao and colleagues used a hybrid pipeline that merged Dlib landmark vectors with CNN features and a support vector machine to classify the four DAiSEE states, reporting 53.42\% overall accuracy on DAiSEE alongside additional evaluation on JAFFE and CK+ \cite{rao_recognition_2021}. Around the same time, Leong compared FaceNet embeddings and landmark-distance streams, each followed by an LSTM, to detect boredom and frustration on DAiSEE \cite{leong_deep_2020}. These two studies showed that educational FER had already moved away from generic basic-emotion labels toward states that matter for learning, but they also showed that single-frame or shallow-temporal architectures struggled with the hard DAiSEE classes.

A clearer step forward came in 2021 with the work of Abedi and Khan, who combined a 2D ResNet-18 with a dilated temporal convolutional network (TCN) and trained it end-to-end on raw video frames. They followed the DAiSEE-defined train, validation and test counts but merged training and validation for fitting, and reported 63.9\% four-class engagement accuracy on the test set, improving over a ResNet+LSTM baseline at 61.15\% \cite{abedi_improving_2021}. Their study also made class imbalance an explicit design issue by using weighted cross-entropy and a balanced-batch sampling strategy. The same temporal direction was followed in educational FER more broadly: Asaju and Vadapalli stacked a ResNet-152 with a BiLSTM and mapped the four DAiSEE states onto a simpler positive--neutral--negative pedagogical scheme, including a small live validation with real students \cite{asaju_adaptive_2022}. Together, these works confirmed that short video clips, rather than isolated frames, were a better unit of analysis for academic affect.

In 2022 the field moved toward attention and transformer-style processing of video. Mehta and colleagues introduced a 3D DenseNet built with spatial, temporal and hybrid non-local self-attention blocks and trained it with class-balanced focal and class-balanced MSE losses \cite{mehta_three-dimensional_2022}. They reported 63.59\% four-class engagement accuracy and 94.78\% binary-engagement accuracy on DAiSEE, where the binary task is obtained by merging the two lowest engagement levels into ``low'' and the two highest into ``high'', and is therefore not directly comparable to the four-level results. They also reported a regression MSE of 0.0347 on DAiSEE engagement and 0.0877 on EmotiW-EP. In the same year, Mandia et al. proposed a Vision Transformer with multi-head self-attention for automatic student engagement estimation on DAiSEE, using imbalance-aware losses and reporting an overall accuracy of about 55.18\%, which the authors describe as outperforming Gupta et al.'s frame-level and video-level baselines by roughly 8--8.78 percentage points \cite{mandia_vision_2022}. The full PDF for this paper was not accessible, so its method details are limited to the available metadata and abstract.

By 2023, two related directions became visible. The first was the hybrid deep-classical pipeline, well represented by Santoni et al., who combined OpenFace-derived facial landmarks, action units, head-pose and gaze features, dimensionality reduction with PCA or SVD, oversampling with SMOTE, and a CNN classifier \cite{santoni_convolutional_2023}. They reported a best maximum accuracy of 77.97\% with the SVD+CNN configuration on a heavily resampled subset of DAiSEE rather than on the original benchmark split, which means the number is not directly comparable with results that follow the standard DAiSEE protocol. This work is important precisely because it shows that DAiSEE research is not only about replacing classical methods with deep networks: it is also about combining deep features with classical preprocessing, feature engineering and imbalance handling. A complementary 2023 study, although not on DAiSEE, illustrates the same hybrid trend in general FER: Mukhiddinov et al.\ used MediaPipe landmarks together with HOG features around those landmarks as input to a custom CNN with batch normalisation and Mish activation, reporting 69.3\% accuracy across seven emotion classes on AffectNet under masked-face conditions \cite{mukhiddinov_masked_2023}. Both studies confirm that handcrafted facial cues remain useful when they are integrated as front-ends to deep models rather than used alone.

The 2024 work continued in three directions. Shiri et al.\ replaced earlier CNN backbones with EfficientNetV2-L and combined it with GRU, Bi-GRU, LSTM or Bi-LSTM heads, reporting their best four-class engagement accuracy of 62.11\% on DAiSEE with the LSTM variant \cite{shiri_detection_2024}; the paper is sometimes cited under a co-author's name but the first author is Shiri. Tian et al.\ proposed a sequential ensemble model evaluated on both DAiSEE and EmotiW-EP, reporting an MSE of 0.0386 on DAiSEE and 0.0610 on EmotiW-EP, with the available abstract claiming stronger performance than prior approaches \cite{tian_predicting_2024}; the full PDF was not accessible, so this entry should be treated as limited verification. Su and colleagues then pushed the field further into the transformer era with a hybrid architecture, called DTransformer, that combines a part-and-sensitive attention network on ResNet-18 patches with a 3D-CNN tube tokenizer feeding a multi-head self-attention transformer, fusing both branches for prediction \cite{su_leveraging_2024}. They reported 64.0\% four-class engagement accuracy on DAiSEE and a regression MSE of 0.0729 on EmotiW-EP, which makes their work one of the strongest current transformer-based engagement detectors.

Two recent 2025 studies show that hybrid deep-classical thinking is still alive. Rianto et al.\ took a graph-based view: they extracted 68 facial landmarks per frame, built a Delaunay triangulation graph, fed it to a graph convolutional network, and combined it with a 1D CNN over 52 action-unit values, training the bimodal model with the LDAM imbalance-aware loss \cite{rianto_engagement_2025}. They followed the DAiSEE engagement-only split and reported a best accuracy of 52.24\% with macro-F1 of 0.49, while honestly noting that minority engagement levels remained largely mispredicted. In general FER, Türkeş and Aydın combined a pre-trained ResNet-50 feature extractor with a Support Vector Machine classifier, reporting 100\% accuracy on the small CK+48 dataset and 100\% on JAFFE in their narrative \cite{turkes_enhanced_2025}. The DAiSEE branch is therefore now bracketed by hybrid graph-and-CNN methods on one side, and by classical hybrid CNN+SVM pipelines on standard FER datasets on the other.

Educational FER has also matured in its target applications. Komaravalli and Janet fine-tuned an Xception-based transfer-learning model on a rebalanced DAiSEE subset and cross-tested it on OL-SFED \cite{komaravalli_detecting_2023}. Gupta and colleagues proposed EDFA, an ensemble that fuses VGG19 and ResNet50 outputs into a rule-based attention index for real-time classroom monitoring, deployed as a web application \cite{gupta_edfa_2023}. Aly combined a ResNet-50 backbone with the Convolutional Block Attention Module and a Temporal Convolutional Network for online student monitoring, reporting throughput of about 28 frames per second on a GPU, 12 on a CPU, and 15 on a Jetson Nano \cite{aly_revolutionizing_2025}, and a follow-up study added a 3D CNN branch with an optimised fusion layer \cite{aly_comprehensive_2025}. Tulsani and colleagues used a hybrid of MTCNN face detection, a Vision Transformer and a BiLSTM inside a learning-management system, with a near-real-time latency of about 0.8 seconds per frame sequence \cite{tulsani_realtime_2025}. Liu and colleagues proposed an adaptive global attention block together with a differential temporal transformer, evaluated on DAiSEE, DFEW and FERV39k, and reported about 11.7 milliseconds per GPU inference \cite{liu_dynamic_2026}. Woolf and colleagues described the Face Readers line of work inside the MathSpring intelligent tutoring system, where attention, head pose and facial cues are used to predict problem-solving outcomes and to trigger real-time interventions \cite{woolf_face_2023}. Together, these examples show that educational FER is no longer only about classification accuracy; it is slowly being connected to the feedback loops of adaptive learning systems.

Taken together, the DAiSEE evidence base evolved through deep-learning baselines, CNN-plus-temporal models, 3D and self-attention architectures, hybrid deep-classical pipelines, EfficientNet-plus-RNN systems, transformer-based detectors and ensemble models, and was complemented by hybrid CNN-and-classical FER work on standard expression datasets. Despite this progress, several limits remain. Datasets for academic emotions are still small and imbalanced; results on the four-level engagement task remain in the low-to-mid sixties even for strong transformer or attention models, and many studies do not use the official DAiSEE split or do not report per-class metrics, which makes head-to-head comparison difficult \cite{abedi_improving_2021,santoni_convolutional_2023,rianto_engagement_2025,shiri_detection_2024,su_leveraging_2024}. Generalisation across cultures and recording conditions remains weak, and visually similar states such as confusion and frustration are still confused \cite{lek_academic_2023,florestiyanto_emotion_2024,vistorte_integrating_2024}. More importantly for this thesis, almost all of the studies above stop at engagement classification, regression or visualisation; very few connect the FER output to a concrete next-task selection or adaptation policy, and live deployment in real classrooms is reported only in a handful of cases. Despite these concerns, facial behaviour remains one of the few signals that can be captured continuously, without attaching sensors to students, without interrupting the lesson, and with hardware that most classrooms and online learners already have. This practical profile is one of the reasons why FER, on its own, continues to be central to recent educational-affect research, and why a webcam-only FER pipeline that is explicitly connected to an adaptive next-task selector still defines an open and useful design point. Table~\ref{tab:fer_comparison} summarises the most relevant FER studies used in this thesis, highlighting their datasets, model designs, target states and practical limitations.



\clearpage
\begin{landscape}
\thispagestyle{empty}
\small
\setlength{\tabcolsep}{3pt}
\renewcommand{\arraystretch}{1.1}
\begin{longtable}{>{\RaggedRight\arraybackslash}p{2.7cm}
                  >{\RaggedRight\arraybackslash}p{1.7cm}
                  >{\RaggedRight\arraybackslash}p{2.3cm}
                  >{\RaggedRight\arraybackslash}p{2.9cm}
                  >{\RaggedRight\arraybackslash}p{1.5cm}
                  >{\RaggedRight\arraybackslash}p{2.0cm}
                  >{\RaggedRight\arraybackslash}p{3.0cm}
                  >{\RaggedRight\arraybackslash}p{3.2cm}}
\caption{Comparison of key FER studies relevant to this thesis}
\label{tab:fer_comparison} \\

\toprule
\textbf{Paper} & \textbf{Dataset} & \textbf{Output labels} & \textbf{Model / Method} & \textbf{Temporal} & \textbf{Main result} & \textbf{Strength / Limitation} & \textbf{Relevance to thesis} \\
\midrule
\endfirsthead

\toprule
\textbf{Paper} & \textbf{Dataset} & \textbf{Output labels} & \textbf{Model / Method} & \textbf{Temporal} & \textbf{Main result} & \textbf{Strength / Limitation} & \textbf{Relevance to thesis} \\
\midrule
\endhead

\bottomrule
\endfoot






\makecell[l]{Asaju and\\ Vadapalli\\ (2021) \cite{asaju_affects_2021}} &
DAiSEE, DISFA+ &
Boredom, Confusion, Frustration, Engagement &
ResNet-50 + Bi-directional LSTM &
Yes &
88\% test accuracy &
Temporal modelling is a strength; DAiSEE split not clearly reported &
Relevant as an early FER-only temporal approach for learning-related emotions \\

\makecell[l]{Asaju and\\ Vadapalli\\ (2022)  \cite{asaju_adaptive_2022}} &
DAiSEE, DISFA+ &
Boredom, Confusion, Frustration, Engagement; mapped to positive / negative / neutral &
ResNet-152 + Bi-LSTM &
Yes &
86\% accuracy on DAiSEE; average F1 = 0.87 &
Strong pedagogical mapping; reasoning layer still limited &
Important because it connects FER output to adaptive feedback logic \\

\makecell[l]{Mehta et al.\\ (2022) \cite{mehta_three-dimensional_2022}} &
DAiSEE, EmotiW-EP &
Boredom, Engagement, Confusion, Frustration &
3D DenseNet-121 + Self-Attention &
Yes &
63.59\% four-class engagement accuracy; 94.78\% binary &
Handles temporal data and label imbalance; some classes remain difficult &
Useful for video-based FER and temporal modelling in education \\

\makecell[l]{Liu et al.\\ (2026) \cite{liu_dynamic_2026}} &
DAiSEE, DFEW, FERV39k &
Boredom, Confusion, Engagement, Frustration &
ResNet18 + Adaptive Global Attention + DTFormer &
Yes &
62.56\% accuracy on DAiSEE &
Efficient inference and advanced temporal modelling; no demographic bias analysis &
Useful as a modern dynamic FER model for subtle academic emotions \\

\makecell[l]{Komaravalli and\\ Janet\\ (2023) \cite{komaravalli_detecting_2023}} &
Customized DAiSEE, OL-SFED &
Boredom, Engagement, Confusion, Frustration &
Xception-based transfer learning &
No &
Frustration 91.87\%, Confusion 83.76\% &
Robust under noisy online conditions; no temporal modelling &
Useful for realistic online learning FER settings \\

\makecell[l]{Rao et al.\\ (2021)\cite{rao_recognition_2021}} &
DAiSEE, JAFFE, CK+ &
Engagement, Frustration, Boredom, Confusion &
Hybrid CNN + landmarks + SVM &
No &
53.42\% accuracy on DAiSEE &
Feature fusion helps, but overall accuracy is low &
Useful as an early FER-only baseline in e-learning \\

\makecell[l]{Leong\\ (2020)\cite{leong_deep_2020}} &
DAiSEE &
Boredom, Frustration &
FaceNet + LSTM / Landmark distances + LSTM &
Yes &
Best reported accuracy: 70.67\% for frustration &
Focus on real-time efficiency; only two categories predicted &
Useful for lightweight FER in adaptive tutoring contexts \\












\makecell[l]{Gupta et al.\\ (2016, arXiv) \cite{gupta_daisee_2016}} &
DAiSEE &
Boredom, Engagement, Confusion, Frustration; 4 ordinal levels &
InceptionNet (frame / video), C3D, LRCN; benchmark &
Yes (LRCN, C3D) &
LRCN top-1 accuracy: Bored. 53.7\%, Eng. 57.9\%, Conf. 72.3\%, Frust. 73.5\% &
Introduces DAiSEE; arXiv preprint; demographic narrow; class imbalance &
Foundational benchmark for educational FER \\

\makecell[l]{Abedi and Khan\\ (2021) \cite{abedi_improving_2021}} &
DAiSEE &
4-class engagement &
ResNet-18 + dilated TCN; weighted cross-entropy &
Yes &
63.9\% four-class engagement accuracy &
arXiv preprint; minority classes still poor; no real-time/F1 reported &
Strong CNN+temporal baseline on the official-style split \\

\makecell[l]{Mandia et al.\\ (2022) \cite{mandia_vision_2022}} &
DAiSEE &
4-class engagement &
Vision Transformer with multi-head attention; imbalance-aware losses &
Yes (transformer over frames) &
$\approx$55.18\% overall accuracy; reportedly +8--8.78 pp over Gupta et al.\ baselines &
Closed access -- limited verification; methodology not fully checked &
Early transformer baseline on DAiSEE \\

\makecell[l]{Santoni et al.\\ (2023) \cite{santoni_convolutional_2023}} &
DAiSEE (custom subset) &
4-class engagement &
OpenFace features (landmarks, AUs, head pose, gaze) + PCA/SVD + SMOTE + CNN &
No (per-frame features, not video clip CNN) &
Best max accuracy 77.97\% (SVD+CNN); avg.\ F1 0.78 &
Custom resampled subset, not the official DAiSEE split; not directly comparable &
Representative hybrid deep-classical pipeline for class imbalance \\

\makecell[l]{Mukhiddinov et al.\\ (2023) \cite{mukhiddinov_masked_2023}} &
AffectNet (masked) &
7 basic emotions &
MediaPipe landmarks + HOG + custom CNN (BatchNorm, Mish) &
No &
69.3\% overall accuracy on AffectNet &
Not DAiSEE; basic emotions; head-pose limitations &
Example of hybrid landmark+HOG+deep FER (general FER context) \\

\makecell[l]{Shiri et al.\\ (2024) \cite{shiri_detection_2024}} &
DAiSEE &
4-class engagement &
EfficientNetV2-L + \{GRU, Bi-GRU, LSTM, Bi-LSTM\} &
Yes (RNN) &
Best 62.11\% accuracy (EfficientNetV2-L + LSTM) &
Per-class metrics not reported; large backbone; no real-time test &
Recent CNN+RNN comparison point \\

\makecell[l]{Tian et al.\\ (2024) \cite{tian_predicting_2024}} &
DAiSEE, EmotiW-EP &
Engagement (regression) &
Sequential ensemble model &
Yes (sequential) &
MSE 0.0386 on DAiSEE; 0.0610 on EmotiW-EP &
Closed access -- limited verification; method details not confirmed &
Ensemble baseline for engagement regression \\

\makecell[l]{Su, He, Luo\\ (2024) \cite{su_leveraging_2024}} &
DAiSEE, EmotiW-EP &
4-class engagement (DAiSEE), regression (EmotiW-EP) &
PANet (ResNet-18 + part attention) + STformer (3D CNN + multi-head transformer) &
Yes (3D + transformer) &
64.0\% accuracy on DAiSEE; MSE 0.0729 on EmotiW-EP &
DAiSEE split not tabulated; not deployed; FER-only &
Strong transformer-based DAiSEE baseline \\

\makecell[l]{Rianto et al.\\ (2025) \cite{rianto_engagement_2025}} &
DAiSEE (engagement only) &
4-class engagement &
68 landmarks $\rightarrow$ Delaunay graph $\rightarrow$ GCN + 52 action units $\rightarrow$ 1D CNN; LDAM loss &
No (per-frame graph + AU streams) &
Best 52.24\% accuracy; macro-F1 0.49 (LDAM) &
Authors note model often predicts only majority classes &
Hybrid landmark/graph + AU + CNN approach to imbalance \\

\makecell[l]{T\"urke\c{s} and Ayd\i n\\ (2025) \cite{turkes_enhanced_2025}} &
CK+48 (also JAFFE narrative) &
6 basic emotions &
Pretrained ResNet-50 features + SVM (RBF) classifier &
No &
100\% accuracy on CK+48; 100\% on JAFFE (narrative) &
Not DAiSEE; small controlled dataset; train/test split not detailed &
Example of hybrid deep+SVM FER (general FER context) \\

  \end{longtable}
\end{landscape}
\clearpage

\section{Adaptive Learning Systems}\label{sec:adaptive_learning_systems}

While the previous section focused on recognising learner states through facial behaviour, adaptive learning systems focus on deciding how instruction should respond to learner needs. These systems are digital tools that change how they teach depending on the individual learner. Instead of presenting the same material in the same order to every student, they adjust things like the content, the pace, the difficulty of the tasks, the feedback given after each answer, hints, and presentation format \cite{strielkowski_span_2025}. A common way to describe this behaviour is as a closed loop: the system collects data from the learner, uses that data to judge how the learner is progressing, and then decides what to show next \cite{gligorea_adaptive_2023}. Reviews of intelligent tutoring systems stress that adaptation is not limited to picking the next exercise: the tutor model may also schedule scaffolding, remediation, motivational messages, interface changes, and assessment pacing \cite{zerkouk_comprehensive_2025}. The goal is to give each learner tasks that match their current level, so that they can keep improving without feeling bored or overwhelmed \cite{zerkouk_comprehensive_2025}.

Early adaptive systems focused on explicit rules written by experts. These were called intelligent tutoring systems (ITS), and examples such as GUIDON and the LISP Tutor were built in the 1980s around hand-crafted if--then rules that decided when to give a hint, when to move on, and when to mark a concept as learned \cite{zerkouk_comprehensive_2025}. The strength of these systems was that their decisions were transparent and easy to explain, but the rules had to be prepared for every topic, which made them hard to extend beyond a single subject. Rule-based logic has not disappeared, though. Sayed et al., for example, used a small set of threshold rules based on a student's penalised grade to move an exercise up or down in difficulty inside an online platform for third-grade mathematics \cite{sayed_ai-based_2023}.

Later work introduced student models that tried to estimate what each learner already knew. Instead of relying only on teacher-written rules, these systems built a picture of the learner from their interaction history and updated it after every answer \cite{zerkouk_comprehensive_2025}. This line of research eventually led to knowledge tracing methods, in which the system predicts how likely the learner is to answer the next question correctly based on their earlier answers. Pardos and colleagues showed that year-long patterns of boredom, engaged concentration, confusion, and gaming-the-system behaviour in ASSISTments logs correlate with end-of-year mathematics test scores \cite{Pardos2014AffectiveStatesStateTests}. Kuling and Zitnik built a recent exercise recommender that tries to pick questions just above the learner's current level \cite{kuling_ken_2025}. Educational psychologists call this range the Zone of Proximal Development: the set of tasks that are challenging but still reachable. Their system ranks candidate questions by predicted difficulty and then chooses items around the 66th percentile, so that the next exercise is hard but not impossible \cite{kuling_ken_2025}.

Over time, researchers began to treat task selection as a decision problem that could be learned from experience. Reinforcement learning (RL) then became a common tool for this. In RL, the system tries different actions, observes the results, and slowly learns which actions lead to better outcomes. Azoulay et al. compared several adaptation methods within the same simulated environment. The study tested different reinforcement learning algorithms, a bandit-style method, and a Bayesian approach that keeps a probability estimate of each learner's ability \cite{azoulay_adaptive_2021}. The Bayesian method reached about 89 to 93 percent of the performance of an ideal algorithm that already knew each student's real ability, while the reinforcement learning methods reached lower but still useful levels \cite{azoulay_adaptive_2021}. The study has become an important reference because it placed very different adaptation families on equal ground, and later work built on its baselines \cite{perez_emotions_2024}.

More recent studies have combined knowledge tracing with reinforcement learning. Fu proposed RL-DKT, a system that first uses a deep knowledge tracing model to estimate the learner's current knowledge state, and then uses a reinforcement learning agent to choose the next task from that estimate \cite{fu_integrating_2025}. Fu tested the method on real learner logs from well-known datasets such as ASSISTments, KDD Cup 2010, and Cognitive Tutor, and reported AUC of 0.874, F1 of 0.812, average completion time of 34.5 minutes, and dropout of 3.2\% \cite{fu_integrating_2025}. Sayed et al. followed a different mix, using a Deep Q-Network to choose how to present the learning content while still using simple rules to move between difficulty levels, and reported learning gains in a small pilot with grade-three mathematics students \cite{sayed_ai-based_2023}. Other studies have used reinforcement learning mainly to decide when and how to give feedback or hints during a session \cite{deshmukh_developing_2025,guzman_developing_2025}. Guzman and Cruz-Mercado, for example, trained a Proximal Policy Optimisation agent that selects problems, hints, and instructional review, and reported normalised learning gains of 0.72 versus 0.58 for a rule-based tutor in a study of ninety introductory Python learners \cite{guzman_developing_2025}. Taken together, these works show that adaptation is increasingly treated as a step-by-step decision problem over tasks, hints, feedback, pacing, or review, in which each chosen action builds on what the learner has done so far.

From 2024, several systems began to close the loop from sensed affect to more than difficulty alone. Pérez et al. used valence and arousal from facial expressions as implicit feedback for a Q-learning difficulty controller \cite{perez_emotions_2024}. Govea and colleagues proposed a deeper variant, in which a Proximal Policy Optimisation agent adapted pace, difficulty and activity type based on a multimodal state that combined face, voice and wrist biometrics \cite{govea_implementation_2024}. Gong described a similar multimodal architecture for e-learning, built around a Deep Q-Network that adjusts difficulty through Item Response Theory, pacing through a Markov-modulated policy and content format between visual, auditory and textual, as well as hints and motivational feedback \cite{gong_design_2025}. Kong and colleagues combined facial expression, voice, heart-rate variability, electrodermal activity and optional electroencephalography into a single joint state of cognitive load and achievement emotions, and a contextual bandit with a Deep Q-Network back end selected interventions such as worked examples, scaffolds, pacing changes or motivational feedback \cite{kong_real-time_2025}. Their randomised evaluation with two hundred and seventy-one learners, in both laboratory and classroom settings, is one of the strongest current examples of a fully working emotion-aware adaptive tutor \cite{kong_real-time_2025}. A scoping review by Fernández-Herrero, also from 2024, confirmed that scaffolding, content adaptation, difficulty, and feedback remain the dominant pedagogical responses when emotion is sensed \cite{fernandez-herrero_evaluating_2024}. These systems are the closest predecessors to an emotion-aware reinforcement-learning tutor, but they either rely on simulation, on static expression corpora, or on sensor stacks that are difficult to deploy in ordinary online classrooms. How these emotion-aware ideas have been put into practice, and how they connect to the adaptive methods described here, is the focus of the next section.

Very recent work extends adaptive tutoring toward generative AI and standardised evaluation. Meng and Yang describe a multi-agent tutor in which Proximal Policy Optimisation selects interventions while neural knowledge tracing and an engagement model track mastery and participation \cite{meng_self-evolving_2025}. Gutierrez, Villegas-Ch and Govea moved to the language-model era with a conversational assistant that adapts feedback, difficulty and communication tone from chat messages \cite{gutierrez_development_2025}. A recent testbed called TutorGym was built to make step-level benchmarking easier and more comparable across studies \cite{weitekamp_tutorgym_2025}.

Evaluating adaptive systems is difficult, because running long classroom studies with many learners is expensive and slow. To deal with this, many studies first test their methods in simulators or on stored learner logs before any live deployment. Azoulay et al. used an artificial environment with 10{,}000 simulated runs to compare adaptation algorithms against a theoretical upper bound \cite{azoulay_adaptive_2021}. Pérez et al. followed a similar style, running 1{,}000 simulations over 200 tasks to compare their approach against earlier reinforcement learning baselines \cite{perez_emotions_2024}. Studies such as those by Fu and by Kuling and Zitnik instead use offline learner datasets and report standard metrics like prediction accuracy and ranking scores on real interaction logs \cite{fu_integrating_2025,kuling_ken_2025}. When live testing is possible, it still tends to be small: Kuling and Zitnik ran a thirteen-week classroom deployment with 38 graduate students \cite{kuling_ken_2025}, and Sayed et al. reported a two-month pilot with 26 grade-three learners \cite{sayed_ai-based_2023}. These practical evaluation approaches matter because they let researchers compare methods and build confidence in them before committing to a full classroom study.

Despite the variety of methods, the main strengths of adaptive learning systems are fairly consistent across the literature. They can keep each learner close to a suitable level of challenge, deliver feedback faster than a human teacher could in a large class, and support thousands of learners at once \cite{gligorea_adaptive_2023}. Large commercial platforms such as Knewton, Carnegie Learning's MATHia, DreamBox, and Smart Sparrow already use these ideas to deliver personalised mathematics and other subjects at scale \cite{strielkowski_span_2025}. Reviews of the field agree that these systems offer clear benefits for engagement and performance when content is well structured and adaptation signals are rich enough \cite{gligorea_adaptive_2023,zerkouk_comprehensive_2025}.

Despite these strengths, the signals that drive the adaptation are still quite narrow. In most of the studies discussed above, the system decides what to show next based only on what the learner did: whether the answer was correct, how long the answer took, how many hints were used, and which items were clicked \cite{azoulay_adaptive_2021,fu_integrating_2025,kuling_ken_2025}. Even when tutors adjust pacing, feedback, or presentation format, the richest multimodal designs still depend on sensors beyond an ordinary webcam \cite{govea_implementation_2024,kong_real-time_2025,gong_design_2025}. How the learner feels during the task, whether they are confused, bored, or frustrated, rarely enters the decision. Azoulay et al. themselves acknowledge this gap, noting that their simulation could check whether the algorithm picked the right difficulty, but could not reflect the learners' satisfaction, motivation, or real learning gain \cite{azoulay_adaptive_2021}. A recent comprehensive review of AI-based intelligent tutoring systems lists affective computing as one of the main open challenges for the field, adding that these systems are still less proficient than human tutors at interpreting students' emotional states \cite{zerkouk_comprehensive_2025}. A recent scoping review reports that more researchers are now trying to add emotion as an additional input channel in intelligent tutoring systems \cite{fernandez-herrero_evaluating_2024}. How these emotion-aware ideas have been put into practice, and how they connect to the adaptive methods described here, is the focus of the next section. To summarise the most relevant adaptive learning studies discussed in this section, Table~\ref{tab:adaptive_systems_comparison} compares their adaptation targets, input signals, decision methods, evaluation settings, main reported results, and relevance to the present thesis.

\clearpage
\begin{landscape}
\thispagestyle{empty}
\small
\setlength{\tabcolsep}{3pt}
\renewcommand{\arraystretch}{1.1}

\begin{longtable}{>{\RaggedRight\arraybackslash}p{2.5cm}
                  >{\RaggedRight\arraybackslash}p{2.1cm}
                  >{\RaggedRight\arraybackslash}p{2.2cm}
                  >{\RaggedRight\arraybackslash}p{2.7cm}
                  >{\RaggedRight\arraybackslash}p{2.0cm}
                  >{\RaggedRight\arraybackslash}p{2.8cm}
                  >{\RaggedRight\arraybackslash}p{3.2cm}
                  >{\RaggedRight\arraybackslash}p{3.4cm}}
\caption{Comparison of key adaptive learning studies relevant to this thesis}
\label{tab:adaptive_systems_comparison} \\

\toprule
\textbf{Paper} & \textbf{Adaptation target} & \textbf{Input signals} & \textbf{Decision method} & \textbf{Evaluation type} & \textbf{Main quantitative results} & \textbf{Main strength / limitation} & \textbf{Relevance to this thesis} \\
\midrule
\endfirsthead

\toprule
\textbf{Paper} & \textbf{Adaptation target} & \textbf{Input signals} & \textbf{Decision method} & \textbf{Evaluation type} & \textbf{Main quantitative results} & \textbf{Main strength / limitation} & \textbf{Relevance to this thesis} \\
\midrule
\endhead

\bottomrule
\endfoot

\makecell[l]{Azoulay et al.\\ (2021)\\ \cite{azoulay_adaptive_2021}} &
Next-task / difficulty selection &
Simulated learner performance, success/failure history &
Comparison of Q-learning, TD, VL, DVRL, Gittins, Rasch+Elo, Bayesian inference &
Simulation &
10{,}000 runs; Bayesian method reached 89--93\% of oracle utility &
Canonical multi-method benchmark; simulation only, no affect or real learners &
Establishes RL/IRT baselines for next-task selection without emotion in state \\

\makecell[l]{Sayed et al.\\ (2023)\\ \cite{sayed_ai-based_2023}} &
Content presentation, exercise navigation, feedback &
Learner model (cognitive, behavioural, affective cues), VARK style &
DQN + threshold rules for difficulty &
Classroom pilot (grade 3, $n=26$, 2 months) &
Post-test gains; largest gains for initially low performers &
Hybrid rules + DQN in deployed e-learning; affect channel not FER-only &
Shows hybrid RL deployment in schools; contrast for full RL + webcam FER policy \\

\makecell[l]{Fu\\ (2025)\\ \cite{fu_integrating_2025}} &
Personalized learning path / next-task sequencing &
Interaction logs, DKT-estimated knowledge state &
RL-DKT: REINFORCE + dynamic knowledge tracing &
Offline datasets (ASSISTments, KDD, Cognitive Tutor) &
AUC 0.874; F1 0.812; avg.\ completion 34.5 min; dropout 3.2\%; path efficiency +12.5\% &
Strong KT+RL sequencing on real logs; no facial affect in state &
Representative performance-only RL sequencing the thesis extends with FER \\

\makecell[l]{Kuling and\\ Zitnik\\ (2025)\\ \cite{kuling_ken_2025}} &
Adaptive exercise / next-question recommendation &
Interaction history, correctness, text embeddings &
KUL-Rec dual-memory Hebbian recommender &
10 offline datasets + classroom ($n=38$, 13 weeks) &
nDCG 0.316 vs 0.265 baseline; Recall@10 0.305 vs 0.211; helpfulness $p<.05$ &
Strong ZPD-targeted recommendation; not RL; no emotion &
Shows next-item adaptation from sparse traces without affect sensing \\

\makecell[l]{Guzman and\\ Cruz-Mercado\\ (2025)\\ \cite{guzman_developing_2025}} &
Problem selection, hints, instructional review &
DKVMN knowledge state; synthetic then real traces &
PPO with multi-faceted reward (performance, efficiency, retention) &
Simulation (10{,}000 synthetic) + human study ($N=90$) &
Normalised learning gains 0.72 vs 0.58 (rule-based) vs 0.45 (static), $p<0.01$ &
End-to-end PPO tutor with human subjects; no facial affect &
Direct PPO comparator for emotion-augmented policy design \\

\makecell[l]{Deshmukh and\\ Sen\\ (2025)\\ \cite{deshmukh_developing_2025}} &
Feedback timing and type (hints, explanations, encouragement) &
Simulated correctness, time on task, hint requests &
Tabular Q-learning &
Simulation vs rule-based and random feedback &
Test accuracy 85\% vs 73\% (rule-based) vs 66\% (random); satisfaction 4.4/5 vs 3.8 &
Simple RL feedback loop; short study; no KT or affect &
Illustrates RL over feedback actions adjacent to thesis intervention set \\

\makecell[l]{P\'erez et al.\\ (2024)\\ \cite{perez_emotions_2024}} &
Difficulty adaptation / next task &
Correctness, BKT mastery, valence/arousal (circumplex) &
Q-learning with emotion as implicit feedback in reward and exploration &
Simulation (1{,}000 runs, 200 tasks) &
Learning gain $\approx$0.49 vs 0.28 (Q-learning) at 200 tasks &
Closest prior: affect in RL loop; simulation and static expression corpus &
Very highly relevant: emotion-aware difficulty RL the thesis generalises to multi-action FER \\

\makecell[l]{Govea et al.\\ (2024)\\ \cite{govea_implementation_2024}} &
Pace, difficulty, activity type, presentation &
Face, voice, wrist biometrics &
CNN--RNN + Proximal Policy Optimization &
Hybrid classroom; multimodal sensing &
Emotion detection 72.4\% $\rightarrow$ 89.3\%; personalization 70.2\% $\rightarrow$ 90.1\% &
Multimodal affect closed loop; adaptation evidence still limited &
Relevant multimodal emotion-to-policy precedent; thesis uses webcam-only sensing \\

\makecell[l]{Gong\\ (2025)\\ \cite{gong_design_2025}} &
Difficulty (IRT), pacing, content format, hints, motivational feedback &
Face (CNN), speech (MFCC+BiLSTM), text (BERT) &
Deep Q-Network adaptive delivery engine &
Controlled e-learning pilot &
Fusion accuracy 0.91; F1 0.90; engagement $\approx$0.92 vs 0.60 baseline &
Multimodal emotion-aware RL over several pedagogical levers; not FER-only &
Very highly relevant multi-action RL design; thesis narrows to facial input \\

\makecell[l]{Kong et al.\\ (2025)\\ \cite{kong_real-time_2025}} &
Worked examples, scaffolds, pacing, motivational feedback, content/interface &
Face, voice, HRV, EDA, optional EEG, interaction logs &
Contextual bandit + Deep Q-Network back end &
RCT ($N=271$); lab + classroom &
Weighted F1 $\approx$0.85; mid-ability post-test gain $\approx$+15 pp; engagement 42\% vs 29\% &
Strongest field evaluation of multi-intervention affect-aware tutor; heavy sensing &
Sets breadth benchmark for intervention set; thesis targets lighter FER+RL stack \\

\makecell[l]{Meng and\\ Yang\\ (2025)\\ \cite{meng_self-evolving_2025}} &
Intervention selection (hints, scaffolding, presentation) &
SAINT+ mastery, behavioural engagement, LLM responses &
PPO multi-agent tutor (RLTS) &
Simulated replay of historical logs &
+28.6\% adaptability; $-$31.2\% recurring errors; $-$24.8\% dropout vs static tutors &
GenAI + KT + RL stack; behavioural not facial engagement &
Contemporary PPO tutor without webcam FER in state \\

\makecell[l]{Weitekamp\\ et al.\\ (2025)\\ \cite{weitekamp_tutorgym_2025}} &
Step-level tutor/student decisions (benchmark) &
ITS step state (223 domains: CTAT, OATutor, etc.) &
Standard API for LLMs, RL, cognitive models &
Benchmark evaluation across existing ITS interfaces &
LLMs near chance on incorrect-action labelling; next-step accuracy $\approx$52--70\% &
Infrastructure for comparable step-level evaluation; not an adaptation algorithm &
Supports simulation-first, step-level evaluation rationale for thesis RL module \\

\end{longtable}
\end{landscape}
\clearpage

\section{Emotion-Aware Adaptation}\label{sec:emotion_aware_adaptation}

While the previous section already introduced recent tutors that combine reinforcement learning with sensed affect, the present section traces how emotion-aware tutoring developed over more than two decades. This line is often called affective tutoring. It began with the question of whether emotion could be recognised during instruction, and only later connected recognition to richer adaptation policies. The narrative below follows that history, from early dialogue-based ideas, through facial and log-based detection, to multimodal reinforcement-learning controllers and conversational assistants based on large language models.

The earliest practical work treated emotion as something that could be detected from conversation with the tutor. AutoTutor, developed by Graesser and his colleagues in the early 2000s, held short natural-language dialogues with students about physics and computer literacy, and its team began to ask whether the same dialogue signals could also show how the learner felt \cite{Graesser2004AutoTutorDialogue}. In 2005, D'Mello and colleagues described an early version of this idea in which affect sensors were combined with the AutoTutor platform to identify boredom, confusion, frustration and engagement from dialogue and simple facial cues \cite{dmello_integrating_2005}. The adaptation at that stage was still limited. The tutor could vary its hints and encouragement, but the main research question was whether emotion could be recognised at all. A few years later, D'Mello and colleagues formalised this direction by showing that ordinary features of tutoring conversation, such as word choice and discourse structure, carried useful information about boredom, confusion, flow and frustration \cite{DMello2008AutomaticDetectionLearnersAffect,dmello_language_2012}. The same line of work also showed that confusion, if handled well, can actually help learning rather than hurt it \cite{dmello2014confusionbeneficial}. Together, these studies established an early view of the field: emotion is visible in how students talk and behave during tutoring, and it is worth recognising, even if the systems themselves still did relatively simple things in response.

A more ambitious direction appeared at the end of the 2000s with the work of Woolf, Burleson, Arroyo, Picard and their colleagues on the Wayang Outpost tutor for school mathematics \cite{woolf_affect-aware_2009}. Instead of relying on dialogue alone, this group built a four-sensor classroom rig that combined a webcam for facial expression analysis, a pressure-sensitive chair for posture, a six-sensor pressure mouse and a wireless skin-conductance bracelet. All four streams fed into machine-learning models that estimated states such as frustration, confidence and engagement, and the tutor used this information to choose the next problem, to offer encouragement, or to bring in a small animated learning companion who mirrored the student's emotion. Around one hundred students used the sensor rig in Fall 2008, and the authors reported that on-target engagement and hint use changed after interventions. This work is historically important because it shows that the core idea of the present thesis, that sensed emotion should influence what the tutor does next, is more than fifteen years old. It also shows how costly early emotion-aware tutoring was: four sensors per student, custom hardware, a dedicated sensor server and a full research team to make it run in a single school.

In the early 2010s, a second strand tried to combine facial expression recognition with text analysis inside a smaller and cheaper affective tutoring system. Lin and colleagues built a prototype for digital art that used an Active Shape Model on webcam images to classify the six basic emotions plus a neutral state, while a semantic voting method inferred learning status from what the student typed \cite{lin_employing_2012}. The system used this information to pick the next teaching strategy and to drive an animated agent that gave sound and animation feedback. The evaluation focused on usability, with System Usability Scale and QUIS surveys, rather than on learning gains. It nevertheless marked an important step, because it showed that a working affective tutor could be built with only a webcam and text input, without the sensor rigs used in the Wayang line.

Around the same time, a third strand focused on making affect detection scale to real classrooms. San Pedro, Baker, Gowda and Heffernan studied the ASSISTments mathematics platform and trained detectors for engaged concentration, boredom and frustration from interaction logs alone, without any camera or physiological sensor \cite{san_pedro_towards_2013}. These detectors were then applied to more than one and a half million actions from over eight thousand students across five school years, and the affect labels were linked to Bayesian knowledge tracing estimates of skill. The authors themselves framed the paper as discovery with models rather than adaptation, and their system did not yet change what the tutor did. The paper is still important for this section for two reasons. First, it showed that large-scale, sensor-free affect detection was possible in real schools. Second, it found that frustration appeared more often both for very weak and for very strong students, a pattern that does not match the simple flow-theory picture often assumed by emotion-driven adaptation rules \cite{san_pedro_towards_2013}. Similar cross-environment work had already warned that boredom in particular has a strong negative link to learning outcomes \cite{Baker2010BetterFrustratedThanBored}, which reinforced the concern that naive adaptation heuristics might not work.

A few years later, Petrovica and Ekenel summarised the state of the field and named the practical trade-off that still shapes it today. In a review of fourteen affective tutoring systems, they argued that the adaptation skills of these systems were still imperfect, mostly because emotion was hard to classify both accurately and unobtrusively \cite{petrovica_emotion_2016}. They contrasted two paths. One was sensor-free, using only log data, which scaled well but lost accuracy. The other was sensor-heavy, using physiological or camera signals, which was more accurate but uncomfortable and expensive. They proposed a middle path, which they called sensor-lite: using only the built-in webcam or microphone that ordinary classrooms already have. This reference is valuable because it states clearly, in 2016, that most affective tutoring systems of that period did not yet have rich real-time adaptation.

As deep learning improved, facial expression recognition became a more serious candidate for the sensor-lite route. Fwa proposed an architectural design for an affective tutoring system for novice programmers, in which the tutor inferred affective states from the webcam and combined them with cognitive signals to decide when to intervene \cite{fwa_architectural_2018}. Work in the same line later studied deeper models for detecting academic emotions such as boredom and frustration from webcam video, focusing on efficient networks that could run in real time \cite{leong_deep_2020}. These studies did not always close the loop into adaptation, but they made it realistic to imagine a tutor that reads the student's face with a single camera and reacts without adding new hardware.

Around 2020, two studies illustrated the two extremes of the trade-off that Petrovica and Ekenel had described. Megahed and Mohammed built an adaptive e-learning system that combined facial expression recognition with a fuzzy-logic rule base \cite{megahed_modeling_2020}. The rules changed the difficulty and the presentation of the content when the detected emotion suggested confusion, boredom or frustration. Because the whole stack used only a webcam and a small set of rules, this design was relatively cheap to deploy. At the other extreme, Standen and colleagues evaluated the MaTHiSiS system for learners with intellectual disabilities in six schools across the United Kingdom, Italy and Spain \cite{standen_evaluation_2020}. Their system combined five sensing channels, including facial expression, eye gaze, body pose, voice and a mobile accelerometer for gesture, and it used this information to select small units of learning content and to change the challenge level. The study used a within-subject design with sixty-seven learners. It found that sessions driven by both affect and achievement produced significantly more engagement and significantly less boredom than sessions driven by achievement alone, but it also found no significant difference in learning gains. This result is one of the clearest pieces of evidence that emotion-aware adaptation can change how students feel without yet changing how much they learn.

From 2024 onward, several systems connected emotion recognition to reinforcement learning, including Pérez et al., Govea and colleagues, Gong, and Kong and colleagues, who were introduced in the previous section \cite{perez_emotions_2024,govea_implementation_2024,gong_design_2025,kong_real-time_2025}. Their designs differ in ways that matter for the present thesis. Pérez, Dapena and Aguilar organised difficulty control around two Q-learning tables, one for performance-based decisions and one for flow-based decisions, and chose when to exploit or explore based on which emotional region the learner occupied in the circumplex model \cite{perez_emotions_2024}. The method was tested in simulation with one thousand runs over two hundred tasks, using expression data rather than a live webcam module. Govea and colleagues reported stronger multimodal detection after training deep reinforcement learning models, but the published evidence for what the adaptation policy changed in hybrid classrooms remained thinner than the detection results \cite{govea_implementation_2024}. Kong and colleagues achieved one of the strongest field evaluations to date, yet noted that mechanism-level findings from their randomised study remained associational rather than strictly causal without tighter experimental control \cite{kong_real-time_2025}. As the scoping review cited above already reported, facial expression remained the dominant sensing channel and scaffolding the most common pedagogical response across twenty-seven affective tutoring systems \cite{fernandez-herrero_evaluating_2024}. A contrasting 2025 direction, taken by Gutierrez, Villegas-Ch and Govea, adapted feedback, difficulty and communication tone from chat messages classified by a fine-tuned RoBERTa model, without using a camera \cite{gutierrez_development_2025}. Together, these lines suggest a split between heavy multimodal reinforcement-learning stacks with broad intervention menus and lighter conversational stacks that trade sensing breadth for deployment ease.

Across all of these phases, several patterns appear again and again. Early work showed that emotion could be detected in tutoring dialogue, but it rarely produced strong adaptation beyond small hint or encouragement changes. Later work showed that emotion could also be detected from facial expressions and from sensor rigs, but the stronger the sensing stack became, the harder it was to deploy in ordinary classrooms. The practical question that Petrovica and Ekenel named almost a decade ago \cite{petrovica_emotion_2016} is therefore still open: how to build an emotion-aware adaptive system that is accurate enough, simple enough to deploy and useful enough pedagogically, all at the same time. The next section returns to this question and sets out the specific research gap that motivates the present thesis.

\section{Research Gap}\label{sec:research_gap}

The previous three sections describe a field that has clearly moved forward. Facial expression recognition in education has grown from handcrafted features and shallow classifiers into transformer-based and lightweight real-time models that fit ordinary webcam pipelines \cite{lek_academic_2023,aly_revolutionizing_2025,aly_comprehensive_2025,liu_dynamic_2026}. Adaptive learning has formalised tutoring as a sequential decision problem and combined knowledge tracing with reinforcement learning to choose exercises, hints, pacing and related interventions from interaction history \cite{azoulay_adaptive_2021,fu_integrating_2025,kuling_ken_2025,axak_adaptive_nodate}. Emotion-aware tutoring has also matured, moving from early dialogue-based prototypes \cite{Graesser2004AutoTutorDialogue,dmello_integrating_2005} and multi-sensor classroom rigs \cite{woolf_affect-aware_2009,standen_evaluation_2020} to recent systems that link emotion to reinforcement learning or to large language models \cite{perez_emotions_2024,govea_implementation_2024,kong_real-time_2025,gong_design_2025,gutierrez_development_2025}. Two recent reviews confirm that both directions are now established research areas rather than isolated experiments \cite{fernandez-herrero_evaluating_2024,zerkouk_comprehensive_2025}. The gap addressed in this thesis is therefore not that emotion-aware adaptation has never been attempted, but that one specific closed-loop configuration---webcam-only affect into a multi-action reinforcement-learning tutor with an evaluable affect channel---is still missing in practice.

When the studies discussed above are read together, several limitations appear at the boundary between facial expression recognition and adaptation. First, most educational facial expression recognition work stops at classification, at an engagement index, or at an instructor dashboard, and does not drive a learnable tutoring policy \cite{rao_recognition_2021,komaravalli_detecting_2023,leong_deep_2020,aly_enhancing_2023,aly_comprehensive_2025,asaju_adaptive_2022}. Second, the strongest reinforcement-learning tutors optimise from cognitive and behavioural traces alone---correctness, timing, hints and mastery estimates---without facial affect in the state \cite{azoulay_adaptive_2021,fu_integrating_2025,kuling_ken_2025,axak_adaptive_nodate}. Comparative work on Proximal Policy Optimisation and Deep Q-Networks in adaptive control illustrates that learnable policies are feasible, but still leaves emotion outside the decision loop \cite{axak_adaptive_nodate}. Third, when emotion-aware systems do close the loop, they tend to follow one of three patterns. They narrow adaptation to a single knob such as difficulty, pacing or tone, rather than a discrete repertoire of pedagogical interventions \cite{perez_emotions_2024,gutierrez_development_2025,megahed_modeling_2020,gong_design_2025}. They depend on multimodal or physiological sensing that is hard to deploy in ordinary online classrooms \cite{govea_implementation_2024,kong_real-time_2025,standen_evaluation_2020}. Or they keep the affect-to-action mapping inside hand-written rules even when other parts of the stack are learned \cite{megahed_modeling_2020,sayed_ai-based_2023}. The closest neighbour to the present work, Pérez and colleagues, does combine facial-expression-derived valence and arousal with Q-learning, but the affect traces come from a static expression corpus rather than from a live webcam module, the controller is organised around difficulty tables rather than a broad intervention set, and the contribution of the emotional channel is not isolated against the same policy without it \cite{perez_emotions_2024}. Although Govea and colleagues and Kong and colleagues report promising deep reinforcement-learning metrics with richer sensing and broader intervention menus, their state representations, reward weights, simulator dynamics and hyperparameters are not always specified in enough detail for independent replication \cite{govea_implementation_2024,kong_real-time_2025}. A field-level scoping review reinforces the pattern that most systems still recognise affect more reliably than they regulate it, with scaffolding the dominant pedagogical response \cite{fernandez-herrero_evaluating_2024}. Even the most carefully evaluated multi-sensor classroom study reported significantly more engagement and less boredom under affect-aware adaptation but no significant difference in learning gains, which suggests that the link from sensed emotion to better learning outcomes still needs more careful tests \cite{standen_evaluation_2020}.

Taken together, these observations point to a precise and defensible gap. The literature reviewed here does not yet contain a system that (i) relies only on a webcam-based facial expression recognition module trained for academic emotions on a multilabel benchmark, following the sensor-lite deployment path that Petrovica and Ekenel argued for classrooms that already have ordinary hardware \cite{petrovica_emotion_2016,aly_comprehensive_2025,liu_dynamic_2026}; (ii) feeds the resulting affective signal into an explicit, learnable reinforcement-learning policy over a \emph{multi-action} pedagogical space---for example hints, scaffolding, pacing, motivation and difficulty adjustment---rather than reducing adaptation to next-item sequencing or a single difficulty rung \cite{perez_emotions_2024,fu_integrating_2025,kuling_ken_2025,kong_real-time_2025}; and (iii) is evaluated under a reproducible simulation protocol that specifies state, reward and transition assumptions, compares emotion-aware and emotion-free variants of the same controller, and reports whether the affect channel changes policy behaviour under matched conditions \cite{azoulay_adaptive_2021,perez_emotions_2024,weitekamp_tutorgym_2025}. Each of these three elements has been studied separately. Webcam-only pipelines are now realistic in real time \cite{aly_comprehensive_2025,liu_dynamic_2026}. Learnable tutoring policies over discrete intervention sets have been demonstrated in simulation and on logs \cite{axak_adaptive_nodate,fu_integrating_2025,kong_real-time_2025}. Reinforcement-learning controllers that take emotion as input have been built, although often under stronger sensing, narrower action spaces, or underspecified training protocols than the configuration targeted here \cite{perez_emotions_2024,govea_implementation_2024,gong_design_2025}. The combination of all three in a single, transparent and scalable closed loop is, however, still missing.

This thesis addresses exactly that combination. The methodology chapter specifies how a lightweight webcam facial expression recognition module, a discrete multi-action reinforcement-learning tutor, and a simulation-based evaluation protocol are built and connected; how affect enters the Markov decision process state (and reward shaping in implementation); and how emotion-aware policies are compared against matched emotion-free baselines under the same student population and hyperparameters. The design follows the sensor-lite direction proposed by Petrovica and Ekenel \cite{petrovica_emotion_2016}: accurate enough to be useful, simple enough to deploy at scale, and explicit enough that decisions can be inspected and replicated. The next chapter describes that architecture, the full MDP specification, and the experimental tests used to answer whether webcam-derived affect improves adaptive control beyond cognitive signals alone.